
% 第 XII 篇，第 7--9 节及参考文献。
\section{在不一定仿射的平滑群预概形上的应用}
\SGARefTarget{sga3:xii:section:7}
\paragraph{\textbf{定理 7.1。}}
\SGARefTarget{sga3:xii:7.1}
\emph{设 \(S\) 是预概形，\(G\) 是在 \(S\) 上平滑、分离且有限型的 \(S\)-群预概形。设 \(G\) 对忠实平坦拟紧拓扑局部地存在极大环面。则：}
\begin{enumerate}
\item[\emph{\SGAInlineRefTarget{sga3:xii:7.1:item:a}(a)}] \emph{映射}
\[
  T\longmapsto \operatorname{Centr}_G(T)
\]
\emph{诱导从 \(G\) 的极大环面集合到 \(G\) 的 Cartan 子群集合的双射。若 \(C\) 对应于 \(T\)，则 \(T\) 是 \(C\) 的唯一极大环面。}
\item[\emph{\SGAInlineRefTarget{sga3:xii:7.1:item:b}(b)}] \emph{设 \(T,T'\) 是 \(G\) 的两个极大环面，\(C,C'\) 为对应的 Cartan 子群。则}
\[
  \operatorname{Transp}_G(T,T')
  =\operatorname{Transp}_G(C,C')
  =\operatorname{Transp}_G(T,C').
\]
\emph{前两项还等同于严格输送子。此外，此函子可由 \(G\) 的一个在 \(S\) 上平滑的闭子预概形表示。环面 \(T,T'\) 以及 Cartan 子群 \(C,C'\) 对平展拓扑局部共轭。}
\item[\emph{\SGAInlineRefTarget{sga3:xii:7.1:item:c}(c)}] \emph{\(G\) 的极大环面及 Cartan 子群对平展拓扑局部地存在。}
\item[\emph{\SGAInlineRefTarget{sga3:xii:theorem:7-1:item:d}\SGAInlineRefTarget{sga3:xii:7.1:item:d}(d)}] \emph{设每个纤维 \(G_s\) 的每个有限子集，都包含于 \(G\) 的一个仿射开集，例如 \(G\) 在 \(S\) 上拟投射，或 \(S\) 是 Artin 的。则 \SGARefLink{sga3:xii:1.10}{1.10} 所定义的 \(G\) 的极大环面函子}
\[
  \mathcal T:(\mathrm{Sch})^\circ/S\longrightarrow(\mathrm{Sets})
\]
\emph{同构于 \(G\) 的 Cartan 子群函子}
\[
  \mathcal C:(\mathrm{Sch})^\circ/S\longrightarrow(\mathrm{Sets})
\]
\emph{且这个 \(G\) 的 Cartan 子群函子可由在 \(S\) 上平滑、分离、有限型的预概形表示。当 \(G\) 在 \(S\) 上拟投射时，此预概形也在 \(S\) 上拟投射；当 \(G\) 在 \(S\) 上仿射，或 \(S\) 是 Artin 的时，它在 \(S\) 上仿射。}
\item[\emph{\SGAInlineRefTarget{sga3:xii:7.1:item:e}(e)}] \emph{设 \(u:G\to G'\) 是 \(S\)-群预概形态射，其中 \(G'\) 在 \(S\) 上平滑、分离且有限型，并假设}
\[
  u_s(G_s)=G'_s
\]
\emph{对每个 \(s\in S\) 成立，即 \(u_s\) 忠实平坦。\footnote{编者未能在第 VI 篇\(_B\) 中找到印本所引的参考；见 M.\ Demazure 和 P.\ Gabriel，《代数群》〔\emph{Groupes algébriques}〕，I，Masson（1970），命题 II.5.1\SGARefLink{sga3:xii:7.1:item:c}{(c)}。}对 \(G\) 的每个极大环面 \(T\)，\(u(T)\) 是 \(G'\) 的极大环面。若 \(G'\) 的纤维连通，则对 \(G\) 的每个 Cartan 子群 \(C\)，\(u(C)\) 是 \(G'\) 的 Cartan 子群；若 \(C\) 是 \(T\) 的中心化子，则 \(u(C)\) 是 \(u(T)\) 的中心化子。在两种情形中，诱导态射 \(T\to u(T)\)（分别为 \(C\to u(C)\)）忠实平坦。}
\item[\emph{\SGAInlineRefTarget{sga3:xii:7.1:item:f}(f)}] \emph{在\textup{\SGARefLink{sga3:xii:7.1:item:e}{(e)}} 的假设下，再设 \(\ker u\) 是 \(G\) 的中心子群。则}
\[
  T\longmapsto u(T)
\]
\emph{是从 \(G\) 的极大环面集合到 \(G'\) 的极大环面集合的双射。若 \(G'\) 的纤维连通，则}
\[
  C\longmapsto u(C)
\]
\emph{同样是从 \(G\) 的 Cartan 子群集合到 \(G'\) 的 Cartan 子群集合的双射；若 \(C'=u(C)\)，则 \(C=u^{-1}(C')\)。}
\end{enumerate}
\paragraph{\textbf{注 7.2。}}
\SGARefTarget{sga3:xii:7.2}
在第 XV 篇 中将看到，即使不对各 \(G_s\) 的有限子集施加上述限制，\SGARefLink{sga3:xii:7.1}{定理 7.1}\textup{\SGARefLink{sga3:xii:7.1:item:d}{(d)}} 的结论仍成立。那里还将证明，\textup{\SGARefLink{sga3:xii:7.1:item:b}{(b)}}、\textup{\SGARefLink{sga3:xii:7.1:item:c}{(c)}}、\textup{\SGARefLink{sga3:xii:7.1:item:d}{(d)}} 和\textup{\SGARefLink{sga3:xii:7.1:item:e}{(e)}} 中仅涉及 Cartan 子群的结论，在更弱的假设下成立：\(G\) 只须对忠实平坦拟紧拓扑局部地存在 Cartan 子群，不必存在极大环面。为此，还需要使用 \(S\) 是 Artin 的情形中的 \SGARefLink{sga3:xii:7.1}{7.1}。当 \(S\) 不是 Artin 的时，第 XV 篇 的方法也将大大简化\textup{\SGARefLink{sga3:xii:7.1:item:c}{(c)}} 和\textup{\SGARefLink{sga3:xii:7.1:item:d}{(d)}} 的证明。
\begin{proof}[\SGARefLink{sga3:xii:7.1}{定理 7.1} 的证明]
\emph{\SGAInlineRefTarget{sga3:xii:7.2:item:a}(a)} 仿照 \SGARefLink{sga3:xii:3.2}{3.2}，可见对 \(G\) 的每个极大环面 \(T\)，
\[
  C=\operatorname{Centr}_G(T)
\]
确实是 \(G\) 的 Cartan 子群，且 \(T\) 由 \(C\) 唯一确定为 \(C\) 的唯一极大环面。只须再说明，\(G\) 的每个 Cartan 子群 \(C\) 都由 \(G\) 的某个极大环面定义，或等价地，\(C\) 存在中心极大环面。
此问题仍对 fpqc 拓扑是局部的，所以可设 \(G\) 存在极大环面 \(T\)。由 \SGARefLink{sga3:xi:theorem:6:2}{XI~6.2}，
\[
  \operatorname{Transp}_G(T,C)
\]
可由 \(G\) 的闭子预概形 \(S'\) 表示，它在 \(S\) 上平滑；此外，由 \SGARefLink{sga3:xii:6.6:item:c}{6.6(c)}，\(S'\to S\) 满。因此，在基变换 \(S'\to S\) 后，可设 \(G\) 有一个截面 \(g\)，使
\[
  \operatorname{ad}(g)\cdot T\subset C.
\]
于是 \(\operatorname{ad}(g)\cdot T\) 是 \(C\) 的极大环面。由 \SGARefLink{sga3:xii:6.6:item:c}{6.6(c)}，它逐纤维地为中心子群，故由 \SGARefLink{sga3:ix:theorem:5:6:item:b}{IX~5.6(b)}，它是中心的。这证明了\textup{\SGARefLink{sga3:xii:7.2:item:a}{(a)}}。
\emph{\SGAInlineRefTarget{sga3:xii:7.2:item:b}(b)} 设 \(T,T'\) 是 \(G\) 的两个极大环面，\(C,C'\) 为对应的 Cartan 子群。以下条件等价：
\[
\begin{array}{ccccc}
(1)\ T\subset T',&
(2)\ T=T',&
(3)\ T\subset C',&
(4)\ C\subset C',&
(5)\ C=C'.
\end{array}
\]
这立即来自\textup{\SGARefLink{sga3:xii:7.2:item:a}{(a)}}。在任意基变换后应用同一结果，得到所断言的输送子与严格输送子之间的等式。
在\textup{\SGARefLink{sga3:xii:7.2:item:a}{(a)}} 中已指出，\(\operatorname{Transp}_G(T,C')\) 可由 \(G\) 的一个在 \(S\) 上平滑的闭子预概形表示，且其结构态射满。由 Hensel 引理，此预概形对平展拓扑局部地有截面。因此，\(T,T'\) 以及 \(C,C'\) 分别对平展拓扑局部共轭。这证明了\textup{\SGARefLink{sga3:xii:7.2:item:b}{(b)}}。
\emph{\SGAInlineRefTarget{sga3:xii:7.2:item:c}(c)} 先设 \(S\) 是局部 Artin 的。若 \(G\) 存在极大环面 \(T\)，刚才证明的共轭定理表明，\(G\) 的极大环面函子 \(\mathcal T\) 可由齐性空间 \(G/N\) 表示，其中 \(N\) 是 \(T\) 在 \(G\) 中的正规化子；由\textup{\SGARefLink{sga3:xii:7.2:item:b}{(b)}}，它平滑。此外，如 VI\(_A\)~\SGARefLink{par:sga3-VIA-3.2.1}{3.2.1} 所指出，\(\mathcal T\) 的每个有限子集都包含于一个仿射开集，因为 \(\mathcal T\) 是 \(G\) 下的齐性空间。
一般而言，\(S\) 的剩余域存在有限扩张 \(k'/k\)，使 \(G_{k'}\) 有极大环面。扩张 \(k'/k\) 由有限平坦基变换 \(S'\to S\) 得到，而由 \SGARefLink{sga3:xi:proposition:2:1:bis}{XI~2.1 bis}，\(G_{k'}\) 的极大环面提升为 \(G_{S'}\) 的极大环面。因此，\(\mathcal T_{S'}\) 可由 \(S'\) 上平滑、分离、有限型的预概形表示，且其每个有限子集都包含于一个仿射开集。其自然下降数据有效。因此 \(\mathcal T\) 可表，由下降可知它在 \(S\) 上平滑、分离且有限型。由平滑性及 Hensel 引理，\(\mathcal T\) 对平展拓扑局部地有截面。这证明了 \(S\) 为 Artin 时的\textup{\SGARefLink{sga3:xii:7.2:item:c}{(c)}} 和\textup{\SGARefLink{sga3:xii:7.1:item:d}{(d)}}。下文将证明，此时 \(\mathcal T\) 实际上在 \(S\) 上仿射。
现在令 \(S\) 任意。为证明\textup{\SGARefLink{sga3:xii:7.2:item:c}{(c)}} 和\textup{\SGARefLink{sga3:xii:7.1:item:d}{(d)}}，这两个断言在 \(S\) 上对 Zariski 拓扑是局部的，可设 \(S\) 仿射，存在与 \(S\) 的所有剩余特征互素的素数 \(\ell\)，且 \(G\) 的纤维的约化秩恒等于 \(r\)。由于假设极大环面对 fpqc 拓扑局部地存在，后一个秩显然局部常值。
选取忠实平坦拟紧态射 \(S'\to S\)，其中 \(S'\) 仿射，且 \(G_{S'}\) 有极大环面 \(T_0\)，令 \(C_0\) 为其中心化子。我断言，存在 \(\ell\) 的适当幂 \(n\)，使
\[
  C_0=\operatorname{Centr}_{G_{S'}}({}_nT_0).
\]
事实上，这立即归约到 \(S'\) 是 Noether 的情形，\SGARefLink{sga3:xi:theorem:6:2}{XI~6.2} 已处理这一情形。固定这样的 \(n\)，置
\[
  M=(\mathbb Z/n\mathbb Z)^r,
  \qquad
  P=\operatorname{Hom}_{S\text{-gr}}(M_S,G).
\]
则 \(P\) 可由 \(({}_nG)^r\) 的一个有限呈示闭子预概形表示，其中
\[
  {}_nG=\operatorname{Hom}_{S\text{-gr}}
    ((\mathbb Z/n\mathbb Z)_S,G).
\]
这也是 \(G\) 上 \(n\) 次幂态射的核，故可由 \(G\) 的有限呈示闭子预概形表示。此外，由 \SGARefLink{sga3:xi:proposition:2:1}{XI~2.1}，\(P\) 在 \(S\) 上平滑。
设 \(s\in S\)。由前文，存在有限可分扩张 \(k''/\kappa(s)\)，使 \(G_{k''}\) 有极大环面 \(T''\)。由于 \({}_nT''\) 在 \(k''\) 上平展，可用 \(k''\) 的有限可分扩张替换它，从而设
\[
  {}_nT''\simeq M_{k''}.
\]
选取平展态射 \(S''\to S\) 以及 \(s\) 上方的一点 \(s''\in S''\)，诱导剩余域扩张 \(\kappa(s'')\simeq k''\)。于是得到 \(P_{S''}\otimes\kappa(s'')\) 在 \(\operatorname{Spec}\kappa(s'')\) 上的截面。由 Hensel 引理，必要时用一个平展邻域替换带点预概形 \((S'',s'')\)，此截面延拓为 \(P_{S''}\) 的截面。等价地，存在同态
\[
  M_{S''}\longrightarrow G_{S''}
\]
诱导给定的同构
\[
  M_{\kappa(s'')}\simeq{}_nT''.
\]
由 \SGARefLink{sga3:ix:theorem:6:4}{IX~6.4}（这里归结为 Nakayama 引理），此同态在 \(s''\) 的一个开邻域上方是闭浸入，可取此邻域为整个 \(S''\)。令 \(H''\) 为其像，\(C''\) 为其在 \(G_{S''}\) 中的中心化子；由 \SGARefLink{sga3:xi:theorem:6:2}{XI~6.2}，\(C''\) 是 \(G_{S''}\) 的闭平滑子预概形。
\paragraph{\textbf{引理 7.3。}}
\SGARefTarget{sga3:xii:7.3}
\emph{在上述关于 \(\ell,n,r\) 的假设下，对 \(S\) 上的每个预概形 \(S''\) 和 \(G_{S''}\) 的每个极大环面 \(T''\)，有}
\[
  \operatorname{Centr}_{G_{S''}}(T'')
  =
  \operatorname{Centr}_{G_{S''}}({}_nT'').
\]
事实上，通过从 \(S'\times_S S''\) 作忠实平坦下降，可归约到 \(G_{S''}\) 存在一个满足上述等式的极大环面 \(T''_1\) 的情形。由\textup{(b)}，\(T''\) 与 \(T''_1\) 对平展拓扑局部共轭，所以同一等式也对 \(T''\) 成立。
在 \(\operatorname{Spec}\kappa(s'')\) 上应用此结果，可知纤维 \(C''_{s''}\) 是 \(G_{s''}\) 的 Cartan 子群。现在，由于 \(G\) 在 \(S\) 上平滑，各纤维 \(G_s\) 的单位连通分支的并是 \(G\) 的开子预概形；这是 EGA IV 中关于平滑态射的一般结果，编者记录了此结果，但未给出确切出处。这个开集显然在群运算下稳定，故为子群，记为 \(G^0\)。它满足与 \(G\) 相同的初始假设，且 \(G\) 的极大环面与 \(G^0\) 的极大环面一一对应。因此，在\textup{(c)} 和\textup{(d)} 的证明中，可以并且将假设 \(G\) 的纤维连通。此时由 \SGARefLink{sga3:xii:6.6:item:a}{6.6(a)}，其 Cartan 子群的纤维也连通。
我断言，\(C''{}^0\) 在 \(s''\) 的一个开邻域上方是 \(G_{S''}\) 的 Cartan 子群。这将完成\textup{\SGARefLink{sga3:xii:7.2:item:c}{(c)}} 的证明。此断言来自下一个引理。
\paragraph{\textbf{引理 7.4。}}
\SGARefTarget{sga3:xii:7.4}
\emph{在 \SGARefLink{sga3:xii:7.1}{7.1} 的假设下，设 \(G\) 的纤维连通。设 \(C\) 是 \(G\) 的在 \(S\) 上平滑的闭群子预概形，\(s\in S\) 满足 \(C_s\) 是 \(G_s\) 的 Cartan 子群。则 \(C^0\) 在 \(s\) 的一个开邻域上是 \(G\) 的 Cartan 子群。}
容易归约到 \(S\) 局部且 \(s\) 为其闭点的情形，并在此情形证明 \(C^0\) 是 \(G\) 的 Cartan 子群。由平坦下降，还可设 \(G\) 存在极大环面 \(T\)。由 \SGARefLink{sga3:xi:theorem:6:2}{XI~6.2}，
\[
  \operatorname{Transp}_G(T,C)
\]
可由 \(G\) 的一个在 \(S\) 上平滑的闭子预概形表示。由关于 \(C_s\) 的假设及共轭\SGARefLink{sga3:xii:6.6:item:a}{定理 6.6(a)}，此输送子在 \(s\) 处的纤维非空。因此，忠实平坦下降将我们归约到输送子在 \(S\) 上有截面的情形，从而归约到 \(C\) 包含 \(G\) 的极大环面 \(T\) 的情形。此时，由 \SGARefLink{sga3:ix:theorem:5:6:item:a}{IX~5.6(a)}，\(T\) 在 \(C^0\) 中是中心的，故
\[
  C^0\subset \operatorname{Centr}_G(T).
\]
这是 \(S\) 上的平滑群预概形之间的包含，它们纤维连通且相对维数相同，即等于它们在 \(s\) 处的共同纤维的维数。因此这是等式，引理得证，\textup{\SGARefLink{sga3:xii:7.2:item:c}{(c)}} 的证明也完成。
\emph{(d)} 保留上述关于 \(\ell,n,r\) 的记号和假设，并假设 \(G\) 的纤维连通。令
\[
  Q:(\mathrm{Sch})^\circ/S\longrightarrow(\mathrm{Sets})
\]
为按 \SGARefLink{sga3:ix:definition:1:4}{IX~1.4} 的意义如下定义的函子
\[
  Q(S')=
  \left\{
  \begin{array}{c}
  \text{\(G_{S'}\) 中的乘法型子群}\\
  \text{其型为 \(M=(\mathbb Z/n\mathbb Z)^r\)}
  \end{array}
  \right\}
\]
规则
\[
  T\longmapsto{}_nT
\]
定义态射
\[
  \varphi:\mathcal T\longrightarrow Q,
\]
由\SGARefLink{sga3:xii:7.3}{引理 7.3}，它是单态射。
我断言，\(Q\) 可由 \(S\) 上分离且有限呈示的预概形表示。如 \SGARefLink{sga3:xi:proposition:3:12:item:a}{XI~3.12(a)} 的证明所指出，有同构
\[
  Q\simeq P_0/\Gamma,
\]
其中 \(P_0\) 是
\[
  P=\operatorname{Hom}_{S\text{-gr}}(M_S,G)
\]
中对应于单态射 \(M_{S'}\to G_{S'}\) 的既开又闭子预概形，如 \SGARefLink{sga3:ix:theorem:6:8}{IX~6.8}，且
\[
  \Gamma=(\operatorname{Aut}_{\mathrm{gr}}(M))_S.
\]
每个纤维 \(G_s\) 的每个有限子集位于一个仿射开集中这一假设，在取闭子预概形及纤维积下稳定。因此，\(({}_nG)^r\)，进而 \(P\)，再进而 \(P_0\)，都继承此性质。故 \(Q\) 可由 \(S\) 上有限呈示的预概形表示。\footnote{编者未具体说明印本对第 V 篇 的引用。}同一论证表明，当 \(G\) 在 \(S\) 上拟投射或仿射时，\(Q\) 也分别在 \(S\) 上拟投射或仿射。在所有情形中，\(Q\) 在 \(S\) 上分离。
\paragraph{\textbf{引理 7.5。}}
\SGARefTarget{sga3:xii:7.5}
\emph{上述态射 \(\varphi:\mathcal T\to Q\) 可由开浸入表示。}
换言之，设 \(H\) 是 \(G\) 中型为 \(M\) 的乘法型子群。必须证明存在开子集 \(U\subset S\)，具有如下性质：对每个 \(S'\to S\)，群 \(H_{S'}\) 形如 \({}_nT_{S'}\)，其中 \(T_{S'}\) 为 \(G_{S'}\) 的某个适当极大环面，当且仅当 \(S'\to S\) 经由 \(U\) 分解。
可设 \(G\) 的纤维的幂零秩恒等于 \(r'\)；由\textup{\SGARefLink{sga3:xii:7.2:item:b}{(b)}}，它局部常值。置
\[
  C=\operatorname{Centr}_G(H),
\]
由 \SGARefLink{sga3:xi:theorem:6:2}{XI~6.2}，这是在 \(S\) 上平滑的闭群子预概形。用 \(C\) 的相对维数为 \(r'\) 的既开又闭子集替换 \(S\)，可设它处处具有这个相对维数。此时，对 \(G\) 的某个极大环面 \(T\) 有 \(H={}_nT\)，当且仅当 \(C^0\) 存在一个处处相对维数为 \(r\) 的中心极大环面 \(T\)，且 \(H={}_nT\)。这给出了子函子 \(U\) 的另一描述。
由平坦下降，可设 \(G\) 存在极大环面 \(T_1\)。考虑子函子
\[
  R=
  \operatorname{Transp}_G
  \bigl(T_1,\operatorname{Centr}_{C^0}(C^0)\bigr)
  \subset
  \operatorname{Transp}_G(T_1,C).
\]
右边的输送子由 \SGARefLink{sga3:xi:theorem:6:2}{XI~6.2} 可由 \(S\) 上的平滑预概形表示。左边的输送子可由诱导的开子预概形表示，这立即来自 \SGARefLink{sga3:ix:theorem:5:6:item:a}{IX~5.6(a)}；特别地，它也在 \(S\) 上平滑。因此，其结构态射是开的。用其像（赋予诱导结构）替换 \(S\)，可设此态射满。将 \SGARefLink{sga3:xii:1.13}{1.13} 用于 \(C^0\)，可知 \(C^0\) 有一个中心极大环面 \(T\)，因为它对 fpqc 拓扑局部地存在这样的环面。其相对维数等于 \(T_1\) 的相对维数，即 \(r\)。关于 \(H\) 的剩余条件为 \(H={}_nT\)，由 \SGARefLink{sga3:ix:corollary:2:10}{IX~2.10}，它又截出 \(S\) 的一个适当的既开又闭子集。引理得证。
\SGARefLink{sga3:xii:7.5}{引理 7.5} 表明，\(\mathcal T\) 可由 \(S\) 上分离且局部有限呈示的预概形表示，实际上它有限呈示。对于后一个断言，可以检查\SGARefLink{sga3:xii:7.5}{引理 7.5} 的证明，看到 \(\varphi\) 是拟紧开浸入；也可以对 fpqc 拓扑局部下降到 \(G\) 有极大环面 \(T\) 的情形，此时
\[
  \mathcal T\simeq G/\operatorname{Norm}_G(T).
\]
这个表达式，或 \SGARefLink{sga3:xi:proposition:2:1}{XI~2.1}，也表明 \(\mathcal T\) 在 \(S\) 上平滑。若 \(G\) 在 \(S\) 上拟投射，则 \(Q\) 与 \(\mathcal T\) 亦然。若 \(G\) 在 \(S\) 上仿射，\(\mathcal T\) 在 \(S\) 上仿射已在 5.4 中证明。这里尚不知道，若不假设 \(G/S\) 仿射，能否选取 \(n\)，使引理 7.5 的开浸入也是闭的。
只须再说明，当 \(S\) 是 Artin 的时，\(\mathcal T\) 在 \(S\) 上仿射。由 EGA I~6.1.7，归约到 \(S\) 是域的谱的情形，并可设该域代数闭。由下文证明的\textup{(f)}，只须对
\[
  G/\operatorname{Centr}(G).
\]
证明断言。由 6.1，此商仿射，故可应用前述仿射情形。这证明了\textup{(d)}。
\emph{\SGAInlineRefTarget{sga3:xii:7.5:item:e}(e)} 由 IX~6.8，\(G'\) 有子环面 \(T'\)，使 \(u\) 诱导忠实平坦态射
\[
  T\longrightarrow T'.
\]
这将 \(T'\) 刻画为 \(G'\) 的子层 \(u(T)\)。令 \(C\) 与 \(C'\) 分别为 \(T\) 与 \(T'\) 的中心化子。我们证明 \(C\to C'\) 平坦，且当 \(G'\) 的纤维连通时忠实平坦。由于 \(C\) 与 \(C'\) 在 \(S\) 上平坦且有限呈示，SGA~1，I~5.9 将我们归约到基为域的情形，并可设该域代数闭。用单位连通分支替换 \(G,G'\) 不改变 \(C^0,C'{}^0\)，所以可设 \(G,G'\) 连通。此时，断言来自 6.6(d) 及\textup{(a)}。
\emph{(f)} 由\textup{(a)} 和\textup{(e)}，只须证明关于 Cartan 子群的断言。\(G\) 的 Cartan 子群 \(C\) 是极大环面的中心化子，故包含 \(G\) 的中心，从而包含 \(\ker u\)。于是
\[
  C=u^{-1}(C'),
  \qquad C'=u(C),
\]
其中 \(C'\) 是\textup{\SGARefLink{sga3:xii:7.5:item:e}{(e)}} 所给的 \(G'\) 的 Cartan 子群。故 \(C\mapsto u(C)\) 是单射。为证明满射性，只须说明，对 \(G'\) 的每个 Cartan 子群 \(C'\)，\(u^{-1}(C')\) 都是 \(G\) 的 Cartan 子群。
此问题对 fpqc 拓扑是局部的。可设 \(G\) 有 Cartan 子群 \(C_1\)。则 \(C'_1=u(C_1)\) 是 \(G'\) 的 Cartan 子群，且由\textup{(b)}，它对 fpqc 拓扑局部地与 \(C'\) 共轭。由于
\[
  u^{-1}(C'_1)=C_1
\]
是 \(G\) 的 Cartan 子群，\(u^{-1}(C')\) 亦然。这证明了\textup{(f)}，也完成了定理的证明。
\end{proof}
或者，为证明 \(u^{-1}(C')\) 是 Cartan 子群，可注意到它在 \(S\) 上平坦，因为 \(u\) 平坦；于是根据定义，SGA~1，I~5.9 将断言归约到基为域的情形，此时 6.6(e) 适用。
\paragraph{\textbf{推论 7.6。}}
\SGARefTarget{sga3:xii:7.6}
\emph{在 7.1\textup{(e)} 的假设下，\(G'\) 对平展拓扑局部地存在极大环面，因而满足对 \(G\) 所施加的初始假设。若再有 \(\ker u\) 为中心的，则 \(G\) 的极大环面和 Cartan 子群函子 \(\mathcal T_G,\mathcal C_G\)，与 \(G'\) 的类似函子 \(\mathcal T_{G'},\mathcal C_{G'}\) 同构。因此，只要这些函子可表，它们就由同构的 \(S\)-预概形表示。}
\paragraph{\textbf{注 7.7。}}
\SGARefTarget{sga3:xii:7.7}
\emph{(a)} 与 \(G\) 在 \(S\) 上仿射的情形不同（实际上只须 \(G\) 的纤维仿射，第 XVI 篇 将说明这一点），\(G\) 的约化秩局部常值，并不蕴含 \(G\) 对 fpqc 拓扑局部地存在极大环面。例如，设 \(S\) 是离散赋值环的谱，\(G_1,G_2\) 是 \(S\) 上平滑、分离、有限型的群概形，使 \(G_1\) 的一般纤维是椭圆曲线，特殊纤维是 \(\mathbb G_m\)，而 \(G_2\) 的一般纤维是 \(\mathbb G_m\)，特殊纤维是 \(\mathbb G_a\)。置
\[
  G=G_1\times_S G_2.
\]
\(G\) 的两个纤维的约化秩均为 \(1\)，但 \(G\) 显然不对 fpqc 拓扑局部地存在极大环面。
另一方面，很有理由相信，对预概形 \(S\) 上平滑、分离、有限型的群，以下条件足以保证对平展拓扑局部地存在极大环面：\(G\) 的纤维的约化秩和阿贝尔秩都是局部常值函数。
\emph{(b)} 在 \SGARefLink{sga3:xii:7.1}{7.1} 的证明中，尤其在\textup{\SGARefLink{sga3:xii:7.1:item:a}{(a)}} 中，我们在未假设 \(S\) 局部 Noether 的情形使用了 \SGARefLink{sga3:xi:theorem:6:2}{XI~6.2}。不过，在所涉及的情形，归约到仿射 Noether \(S\) 是直接的。
以下是 \SGARefLink{sga3:xii:7.1:item:b}{7.1(b)} 的一个变体。
\paragraph{\textbf{命题 7.8。}}
\SGARefTarget{sga3:xii:7.8}
\emph{设 \(G\) 是纤维连通的平滑有限呈示 \(S\)-群预概形，\(H\) 是平滑 \(S\)-群预概形，}
\[
  i:H\longrightarrow G
\]
\emph{是 \(S\)-群单态射，将 \(H\) 作为 \(G\) 的 \(S\)-子群。对 \(G\) 的每个 Cartan 子群 \(C\)，\footnote{证明只要求 \(C\) 是 \(G\) 的平滑子群，其几何纤维是极大环面的子群的连通中心化子。}输送子 \(\operatorname{Transp}_G(C,H)\) 可由 \(G\) 的一个在 \(S\) 上平滑的闭子预概形表示。}
输送子可由 \(G\) 的有限呈示闭子预概形表示，这来自 \SGARefLink{sga3:xi:corollary:6:11}{XI~6.11}，因为由 \SGARefLink{sga3:xii:6.6:item:b}{6.6(b)}，当 \(G\) 的纤维连通时，\(C\) 的纤维也连通。为证明平滑性，标准过程先归约到仿射 Noether \(S\)，再到局部 Artin \(S\)，最后通过下降归约到剩余域代数闭的情形。此时，\(C\) 有极大环面 \(T\)，它也是 \(G\) 的极大环面。除非纤维 \(H_0\) 的约化秩和幂零秩分别等于 \(G_0\) 的相应秩，否则输送子为空。在这些秩相等的条件下，利用 \(C\) 的连通性以及 \(H\) 的极大环面在 \(H\) 中的中心化子的平滑性，可见
\[
  \operatorname{Transp}_G(C,H)
  =
  \operatorname{Transp}_G(T,H).
\]
右边由 \SGARefLink{sga3:xi:proposition:2:5}{XI~2.5} 平滑，故左边也平滑。
同一论证证明下文更一般的断言 \SGARefLink{sga3:xii:7.9:item:b:02}{7.9(b)}。
\paragraph{\textbf{命题 7.9。}}
\SGARefTarget{sga3:xii:7.9}
\emph{设 \(G\) 和 \(i:H\to G\) 如 \SGARefLink{sga3:xii:7.8}{7.8}。再假设对每个 \(s\in S\)，\(H_s\) 连通，且与 \(G_s\) 具有相同的约化秩和幂零秩；等价地，\(H_s\) 包含 \(G_s\) 的一个 Cartan 子群。则：}
\begin{enumerate}
\item[\emph{(a)}] \emph{\(\operatorname{Norm}_G(H)\) 可由 \(G\) 的一个在 \(S\) 上有限呈示的闭子预概形表示，且典范单态射}
\[
  H\longrightarrow\operatorname{Norm}_G(H)
\]
\emph{是开浸入。因此 \(i\) 是浸入，且}
\[
  H=\operatorname{Norm}_G(H)^0.
\]
\item[\emph{(b)}] \emph{对 \(G\) 的每个 Cartan 子群 \(C\)，\(\operatorname{Transp}_G(C,H)\) 是 \(G\) 的闭子预概形，在 \(S\) 上平滑且结构态射满。若 \(C\) 是 \(G\) 的极大环面 \(T\) 的中心化子，则还有}
\[
  \operatorname{Transp}_G(T,H)
  =
  \operatorname{Transp}_G(C,H).
\]
\item[\emph{(c)}] \emph{设 \(C\) 是 \(H\) 的群子预概形。则 \(C\) 是 \(H\) 的 Cartan 子群，当且仅当它是 \(G\) 的 Cartan 子群。}
\item[\emph{(d)}] \emph{若 \(G\) 对平展拓扑或 fpqc 拓扑局部地存在 Cartan 子群，或分别存在极大环面，则 \(H\) 亦然。}
\end{enumerate}
\begin{proof}
\emph{\SGAInlineRefTarget{sga3:xii:7.9:item:a:02}(a)} \(\operatorname{Norm}_G(H)\) 可由 \(G\) 的一个在 \(S\) 上有限呈示的闭子预概形表示，这包含在 \SGARefLink{sga3:xi:corollary:6:11}{XI~6.11} 中。由于 \(H\) 平滑，因而在 \(S\) 上平坦且局部有限呈示，VI\(_B\)~\SGARefLink{cor:sga3-VIB-2.6}{2.6} 将
\[
  H\longrightarrow\operatorname{Norm}_G(H)
\]
是开浸入的断言归约到纤维。因此可归约到 \(S=\operatorname{Spec}k\)，其中 \(k\) 代数闭。
此时断言归结为说明，对应的 Lie 代数映射是同构，\footnote{编者未能在第 VI 篇\(_B\) 中找到印本所引的参考；见 M.\ Demazure 和 P.\ Gabriel，《代数群》〔\emph{Groupes algébriques}〕，I，Masson（1970），推论 II.5.6。}或等价地，
\[
  (\mathfrak g/\mathfrak h)^H=0,
\]
其中 \(\mathfrak g,\mathfrak h\) 是 \(G,H\) 的 Lie 代数，上标表示不变量；比较 II~5.2.3(i)。由假设，\(H\) 包含 \(G\) 的一个 Cartan 子群 \(C\)，即 \(G\) 的某个极大环面 \(T\) 的中心化子。因此，只须证明
\[
  (\mathfrak g/\mathfrak h)^T=0.
\]
由 \(T\) 的表示的完全可约性（\SGARefLink{sga3:i:proposition:4-7-3}{I~4.7.3}），这来自
\[
  (\mathfrak g/\mathfrak c)^T=0,
  \qquad \mathfrak c=\operatorname{Lie}(C).
\]
最后这个等式等价于
\[
  \mathfrak g^T=\mathfrak c.
\]
它表示中心化子 \(C\) 与 \(T\) 的正规化子具有相同的 Lie 代数，这来自 \(C\) 是该正规化子的开子群这一事实（\SGARefLink{sga3:xi:proposition:5:9}{XI~5.9}）。这证明了\textup{\SGARefLink{sga3:xii:7.9:item:a:02}{(a)}}。
\emph{\SGAInlineRefTarget{sga3:xii:7.9:item:b:02}(b)} 已在 \SGARefLink{sga3:xii:7.8}{7.8} 中证明。
\emph{\SGAInlineRefTarget{sga3:xii:7.9:item:c:02}(c)} 由于 \(H\) 是 \(G\) 的子预概形，断言立即归约到代数闭基域，此时它直接来自关于 \(H\) 的假设。
\emph{\SGAInlineRefTarget{sga3:xii:7.9:item:d:02}(d)} 来自\textup{\SGARefLink{sga3:xii:7.9:item:b:02}{(b)}}、\textup{\SGARefLink{sga3:xii:7.9:item:c:02}{(c)}} 以及 Hensel 引理 \SGARefLink{sga3:xi:proposition:1:10}{XI~1.10}。
\end{proof}
\paragraph{\textbf{推论 7.10。}}
\SGARefTarget{sga3:xii:7.10}
\emph{设 \(G,H\) 如 \SGARefLink{sga3:xii:7.9}{7.9}，并设对 \(S\) 上的每个代数闭域 \(k\)，\(G_k(k)\) 中正规化 \(H_k\) 的每个元素都属于 \(H_k(k)\)。则 \(H\) 是 \(G\) 的闭子预概形，且等于自身的正规化子。}
这立即来自 \SGARefLink{sga3:xii:7.9:item:a:02}{7.9(a)}。我们特别将其用于 \(G\) 的 Borel 子群，更一般地用于抛物子群。
\paragraph{\textbf{推论 7.11。}}
\SGARefTarget{sga3:xii:7.11}
\emph{设 \(G,H\) 如 \SGARefLink{sga3:xii:7.9}{7.9}。则 \(H\) 包含 \(G\) 的每个在 \(G\) 中为中心、在 \(S\) 上平坦且有限呈示的群子预概形 \(Z\)。}
照例先归约到仿射 Noether \(S\)，再到 Artin \(S\)；此时 \SGARefLink{sga3:xii:7.1}{7.1} 的假设成立。由 \SGARefLink{sga3:xii:7.9:item:c:02}{7.9(c)}，可设 \(H\) 包含 \(G\) 的 Cartan 子群 \(C\)，只须再证 \(C\supset Z\)。由于 \(S\) 是 Artin 的，商
\[
  G'=G/Z
\]
可由 \(S\) 上有限型群预概形表示；典范态射 \(u:G\to G'\) 忠实平坦，且由 VI\(_A\)~\SGARefLink{sga3:xii:3.2}{3.2}，其核为 \(Z\)。群 \(G'\) 在 \(S\) 上平滑。因此，\SGARefLink{sga3:xii:7.1:item:f}{定理 7.1(f)} 表明，对 \(G'\) 的某个 Cartan 子群 \(C'\)，有 \(C=u^{-1}(C')\)，故 \(C\) 包含 \(Z\)。
\paragraph{\textbf{推论 7.12。}}
\SGARefTarget{sga3:xii:7.12}
\emph{设 \(G,G'\) 是平滑有限呈示 \(S\)-群预概形，\(u:G\to G'\) 是忠实平坦群同态，即对每个 \(s\in S\)，\(u_s\) 忠实平坦。设 \(\ker u\) 是中心的，且 \(G\) 的纤维连通。则}
\[
  H'\longmapsto H=u^{-1}(H')
\]
\emph{给出以下两者之间的双射：}
\begin{itemize}
\item \emph{\(G'\) 的在 \(S\) 上平滑且有限呈示的群子预概形 \(H'\)，在每个 \(s\in S\) 处与 \(G'\) 具有相同的约化秩和幂零秩；以及}
\item \emph{\(G\) 的在 \(S\) 上平滑且有限呈示的群子预概形 \(H\)，在每个 \(s\in S\) 处与 \(G\) 具有相同的约化秩和幂零秩。}
\end{itemize}
\emph{此外，\(H\) 的纤维连通，当且仅当 \(H'\) 的纤维连通。}
设 \(H\) 是 \(G\) 的具有上述性质的群子预概形。由 \SGARefLink{sga3:xii:7.11}{7.11}，它包含 \(Z=\ker u\)，故忠实平坦下降给出
\[
  H=u^{-1}(H')
\]
其中 \(H'\) 是唯一确定的 \(G'\) 的群子预概形。利用 \SGARefLink{sga3:xii:6.6:item:e}{6.6(e)}，立即可验证 \(H'\) 具有所断言的性质。若 \(H\) 的纤维连通，则 \(H'=H/Z\) 亦然。
反之，取具有上述性质的 \(H'\subset G'\)，置 \(H=u^{-1}(H')\)。由 \SGARefLink{sga3:xii:6.6:item:e}{6.6(e)}，只须再说明 \(H\) 在 \(S\) 上平滑，且在 \(H'\) 的纤维连通时其纤维也连通。由于 \(u\) 平坦，且 \(H'\) 在 \(S\) 上平坦，\(H\) 已在 \(S\) 上平坦。因此只须在几何纤维上检验平滑性，或分别检验连通性，从而归约到代数闭基域。
此时 \(H'\) 包含 \(G'\) 的 Cartan 子群 \(C'\)，而 \(H\) 包含其逆像 \(C\)。由 \SGARefLink{sga3:xii:6.6:item:e}{6.6(e)}，\(C\) 是 \(G\) 的 Cartan 子群，故平滑且连通。现在结论来自下一个引理。
\paragraph{\textbf{引理 7.13。}}
\SGARefTarget{sga3:xii:7.13}
\emph{设 \(G,G'\) 是 \(S\) 上平坦且有限呈示的群预概形，\(u:G\to G'\) 是忠实平坦群同态，\(C'\) 是 \(G'\) 的在 \(S\) 上有限呈示的群子预概形，使 \(C=u^{-1}(C')\) 在 \(S\) 上平滑，或分别具有连通纤维。则对 \(G'\) 的每个在 \(S\) 上有限呈示、包含 \(C'\)、且在 \(S\) 上平滑（或分别具有连通纤维）的群子预概形 \(H'\)，其逆像 \(H=u^{-1}(H')\) 在 \(S\) 上平滑，或分别具有连通纤维。}
如前所述，这立即归约到 \(S\) 是域的谱的情形。由 VI\(_B\)~9，群 \(H\) 是 \(H/C\) 上的主 \(C\)-丛，而由第 IV 篇，
\[
  H/C\simeq H'/C'
\]
由于 \(H'\) 平滑（分别为连通），\(H'/C'\) 也平滑（分别为连通），从而 \(H/C\) 亦然。由假设，\(C\) 具有同一性质，故丛 \(H\) 也具有此性质。这证明了引理及\SGARefLink{sga3:xii:7.12}{推论 7.12}。
\section{半单元素；不一定仿射的群概形中极大环面的并与交}
\SGARefTarget{sga3:xii:section:8}
本节中，\(G\) 始终表示在 \(S\) 上平滑且有限呈示、纤维连通的 \(S\)-群预概形。
先设 \(S\) 是代数闭域 \(k\) 的谱。当 \(G\) 仿射时，\(G(k)\) 的半单元素的概念已在\emph{BIBLE}~4，第 4 节定义；立即可验证，此概念在基域的代数闭扩张 \(k'/k\) 下不变。此外，由\emph{BIBLE}~6，定理 5(c)，元素 \(g\in G(k)\) 半单，当且仅当它包含于 \(G\) 的一个极大环面。
不再假设 \(G\) 仿射时，Chevalley 定理给出 \(G\) 作为阿贝尔簇按平滑连通仿射代数群的扩张的典范表达式：
\begin{equation}
\tag{*}\label{sga3:xii:lemma:7-13:equation:star}
  e\longrightarrow V\longrightarrow G\longrightarrow A
  \longrightarrow e.
\end{equation}
若 \(g\in G(k)\) 是 \(V(k)\) 的半单元素，则称其为\emph{半单}的。由于 \(V\) 的极大环面恰好是 \(G\) 的极大环面，这等价于说 \(g\) 属于 \(G\) 的一个极大环面。此概念在基域的代数闭扩张下仍不变。
再设 \(S=\operatorname{Spec}k\)，其中 \(k\) 是任意域，并取 \(g\in G\)。选取 \(\kappa(g)\) 的代数闭扩张 \(K\) 后，点 \(g\) 是 \(G\) 的一个在 \(k\) 的代数闭扩张 \(K\) 中取值的几何点 \(g'\) 的像。若 \(g'\) 半单，则称 \(g\) 半单。前文表明，这不依赖于 \(g'\) 的选择。若 \(k'/k\) 是任意扩张，则 \(G_{k'}\) 的半单元素集合是 \(G\) 的半单元素集合 \(G^{\mathrm{ss}}\) 的逆像。
最后，对任意 \(S\)，若点 \(g\in G\) 在其纤维 \(G_s\) 中半单，则称它半单。因此，对每个态射 \(S'\to S\)，子集 \(G^{\mathrm{ss}}_{S'}\) 是 \(G^{\mathrm{ss}}\) 的逆像。假设 \SGARefLink{sga3:xii:1.10}{1.10} 定义的极大环面函子 \(\mathcal T\) 可由 \(S\) 上有限呈示的预概形表示。例如，若 \(G\) 对 fpqc 拓扑局部地存在极大环面，由 \SGARefLink{sga3:xii:7.1:item:d}{7.1(d)}，至少当 \(G\) 在 \(S\) 上拟投射时如此。令 \(T\) 为 \(G_{\mathcal T}\) 的典范极大环面，即“\(G\) 的万有极大环面”，考虑由投影 \(G_{\mathcal T}\to G\) 诱导的态射
\[
  u:T\longrightarrow G
\]
由定义立即可知，\(G^{\mathrm{ss}}\) 是此态射的像。得到：
\paragraph{\textbf{命题 8.1。}}
\SGARefTarget{sga3:xii:8.1}
\emph{\(G\) 的半单元素集合 \(G^{\mathrm{ss}}\) 局部可构造。因此，当 \(S\)，从而 \(G\)，拟紧且拟分离时，它可构造。}
照例归约到仿射 Noether \(S\)。由可构造性的标准 Noether 判据（EGA \(0_{\mathrm{III}}\)，9.2.3），还可设 \(S\) 整，只须说明存在非空开集 \(U\subset S\)，使 \(G^{\mathrm{ss}}|_U\) 可构造。缩小 \(U\)，必要时用有限覆盖替换它，可设 \(G\) 在 \(S\) 上分离且包含极大环面。由 7.1(d)，\(\mathcal T\) 此时可由 \(S\) 上有限呈示的预概形表示；其万有环面也如此，故它在 \(G\) 中的像可构造。
再设 \(k\) 是域，令 \(H\) 为由上述态射生成的 \(G\) 的子群（VI\(_B\)~\SGARefLink{sga3:xii:1.2}{1.2}）。它是 \(G\) 的平滑群子概形，且因万有环面连通而连通；其构造与每个基域扩张交换，\(\mathcal T\) 的构造也如此。
当 \(k\) 代数闭时，\(H\) 也是由 \(G\) 的极大环面，或等价地由 \(G\) 的所有环面生成的 \(G\) 的代数子群。它同样是包含 \(G(k)\) 的所有半单元素的 \(G\) 的最小代数子群。只要 \(k\) 无限，这些刻画仍然有效，因为第 XIV 篇 将证明，\(\mathcal T\) 的 \(k\)-有理点集合在 \(\mathcal T\) 中稠密。
子群 \(H\) 在 \(G\) 中不变。事实上，只须在取代数闭域后检验此事；此时 \(H\) 与 \(G\) 平滑，只须检验在内自同构 \(\operatorname{int}(g)\)（\(g\in G(k)\)）下的稳定性，这显然成立。甚至可以说明 \(H\) 在 \(G\) 中是特征子群，即在 \(\operatorname{Aut}_{k\text{-}\mathrm{gr}}(G)\) 下稳定。
由 \SGARefLink{sga3:xii:6.6:item:d}{6.6(d)} 可知 \(G/H\) 的约化秩为零。更确切地，若 \(K\) 是 \(G\) 的不变代数子群，则在代数闭域上，\(G/K\) 的极大环面是 \(G\) 的极大环面的直接像。因此，\(G/K\) 的约化秩为零，当且仅当 \(K\) 包含 \(G\) 的所有极大环面；或在任意域上，当且仅当 \(K\supset H\)。因此，\(H\) 是使 \(G/H\) 的约化秩为零的 \(G\) 的最小不变代数子群。它也是与 \(G\) 具有相同约化秩的 \(G\) 的最小代数子群。
最后，\(H\) 仿射。这仍可在取代数闭包后检验。在正合列 \((*)\) 中，\(G\) 的每个极大环面都包含于 \(V\)；故 \(H\subset V\)，而 \(V\) 仿射。我们证明了：
\paragraph{\textbf{命题 8.2。}}
\SGARefTarget{sga3:xii:8.2}
\emph{设 \(G\) 是域 \(k\) 上的平滑连通代数群，\(\bar k\) 是 \(k\) 的代数闭包。存在 \(G\) 的代数子群 \(H\)，使 \(H_{\bar k}\) 是 \(G_{\bar k}\) 中由其极大环面，或等价地由 \(G_{\bar k}\) 的所有环面生成的代数子群。群 \(H\) 也可刻画为与 \(G\) 具有相同约化秩的 \(G\) 的最小代数子群，或使 \(G/H\) 的约化秩为零的 \(G\) 的最小不变代数子群。它是 \(G\) 的平滑、连通、不变、仿射子群，且其构造与每个基域扩张交换。}
用 \(G/H\) 给出的刻画有如下推论。
\paragraph{\textbf{推论 8.3。}}
\SGARefTarget{sga3:xii:8.3}
\emph{设 \(G\) 是代数闭域 \(k\) 上的平滑连通代数群。则 \(G\) 的约化秩为零——等价地，按 \SGARefLink{sga3:xii:8.2}{8.2} 的记号，\(H=(e)\)——当且仅当 \(G\) 是阿贝尔簇按平滑连通幺幂代数群的扩张，后者即由同构于加法群 \(\mathbb G_a\) 的群逐次扩张得到的群。}
事实上，Chevalley 正合列 \((*)\) 将断言归结为：平滑连通仿射群 \(V\) 的约化秩为零，当且仅当它幺幂；这包含在\emph{BIBLE}~\SGARefLink{sga3:xii:6.4}{6.4}，定理 4，\SGARefLink{sga3:xii:bibliography}{推论 3} 中。因此，按 \SGARefLink{sga3:xii:8.2}{8.2} 的记号，\(H\) 是 \(G\) 的最小不变代数子群，使 \((G/H)_{\bar k}\) 是阿贝尔簇按平滑连通仿射幺幂代数群的扩张。
\paragraph{\textbf{推论 8.4。}}
\SGARefTarget{sga3:xii:8.4}
\emph{按 \SGARefLink{sga3:xii:8.2}{8.2} 的记号，\(G=H\)，或等价地 \(G_{\bar k}\) 由其极大环面生成，当且仅当 \(G\) 仿射，且每个同态 \(G_{\bar k}\to\mathbb G_a\) 都平凡。}
\paragraph{\textbf{注 8.5。}}
\SGARefTarget{sga3:xii:8.5}
\begin{enumerate}
\item[\emph{\SGAInlineRefTarget{sga3:xii:8.5:item:a}(a)}] 设 \(V\) 是 \(G_{\bar k}\) 的最大平滑连通仿射代数子群，使 \(G_{\bar k}\) 是阿贝尔簇按 \(V\) 的扩张。由 Rosenlicht 的结果，\footnote{编者要求在此补充参考文献。}众所周知，当 \(k\) 不完美时，\(V\) 一般不在 \(k\) 上定义：未必存在 \(G\) 的代数子群 \(V_0\)，使 \((V_0)_{\bar k}=V\)。不过，若 \(V\) 由其极大环面生成——等价地，\(V\) 不存在同构于 \(\mathbb G_{a,\bar k}\) 的商——则确实存在这样的群，即 \SGARefLink{sga3:xii:8.2}{8.2} 的群 \(H\)。因此，在这个有理性问题以及许多其他问题中（比较第 XIV 篇，第 6 节），所有困难都来自幺幂群，因而来自加法群；而足够多的乘法型群的存在，则使下降能够进行。
\item[\emph{(b)}] 按 \SGARefLink{sga3:xii:8.2}{8.2} 的记号，令 \(H'\) 为 \(G/H\) 的交换子群在 \(G\) 中的逆像。则 \(H'\) 是 \(G\) 的最小不变代数子群，使 \((G/H')_{\bar k}\) 是阿贝尔簇按平滑连通仿射幺幂代数群的交换扩张。它仍是 \(G\) 的平滑连通仿射子群。
设特征 \(p\) 为正，令 \(H''\) 为充分大的 \(n\) 所对应的 \(p^n(G/H')\) 在 \(G\) 中的逆像。则 \(H''/H'\) 是阿贝尔簇，且 \(H''\) 是 \(G\) 的最小不变代数子群，使 \(G/H''\) 是交换平滑连通仿射幺幂代数群。可知，按\textup{\SGARefLink{sga3:xii:8.5:item:a}{(a)}} 的记号，\(V=H_{\bar k}\)，当且仅当 \(H''=G\)，又当且仅当每个同态 \(G_{\bar k}\to\mathbb G_a\) 都平凡。
\end{enumerate}
现在按照 \SGARefLink{sga3:xii:4.10}{4.10} 的思路，将\SGARefLink{sga3:xii:section:4}{第 4 节}发展的约化中心概念推广到纤维连通的平滑群。令 \(\mathcal Z\) 为如下定义的 \(G\) 的子函子
\[
\begin{split}
  \mathcal Z(S')=\{\, &g'\in G_{S'}(S'):\text{ 对每个 }
  S''\to S'\text{ 以及每个极大环面 }T''\subset G_{S''},\\
  &g'\text{ 的拉回 }g''\text{ 是 }
  T''\text{ 在 }S''\text{ 上的截面}\,\}.
\end{split}
\]
令 \(\mathcal T(S')\) 为 \(G_{S'}\) 的极大环面集合，\(\widetilde{\mathcal T}(S')\) 为对 \((T',g')\) 的集合，其中 \(T'\) 是 \(G_{S'}\) 的极大环面，\(g'\) 是 \(T'\) 在 \(S'\) 上的截面。则 \(\widetilde{\mathcal T}\) 是
\[
  \mathcal T\times_S G=\mathcal T_G,
\]
的子函子，且用此记号，还可写为
\[
  \mathcal Z=
  \prod_{\mathcal T_G/G}
  \widetilde{\mathcal T}/\mathcal T_G.
\]
利用 \SGARefLink{sga3:xi:theorem:6:8}{XI~6.8} 和 \SGARefLink{sga3:xii:7.1:item:d}{7.1(d)}（后者在适当假设下将 \(\mathcal T\) 由平滑 \(S\)-预概形表示），可以在相应假设下推出 \(\mathcal Z\) 的可表性。我们将在下文更直接地得到它。
\paragraph{\textbf{定义 8.6。}}
\SGARefTarget{sga3:xii:8.6}
\emph{设 \(G\) 是平滑有限呈示且纤维连通的 \(S\)-群预概形。若上述函子 \(\mathcal Z\)（显然是 \(G\) 的子群）可由乘法型群表示，则称 \(G\) 存在\emph{约化中心}。此时，称 \(\mathcal Z\) 为 \(G\) 的约化中心。}
若 \(Z\) 是 \(G\) 的约化中心，则对每个基变换 \(S'\to S\)，\(Z_{S'}\) 是 \(G_{S'}\) 的约化中心。因此，约化中心的存在性对 fpqc 拓扑是局部的。这个术语是恰当的，因为 \(\mathcal Z\) 在
\[
  \operatorname{Aut}_S(G),
\]
下不变，因而是 \(G\) 的不变子群；由 \SGARefLink{sga3:ix:corollary:5:5}{IX~5.5}，它是中心的。
这里用以下陈述替代\SGARefLink{sga3:xii:4.5}{引理 4.5}。
\paragraph{\textbf{引理 8.7。}}
\SGARefTarget{sga3:xii:8.7}
\emph{设 \(u:H\to G\) 是中心同态，其中 \(H\) 是 \(S\) 上有限型乘法型群。假设对 \(S\) 上的每个代数闭域 \(k\)，子群 \(u_k(H_k)\) 包含于 \(G_k\) 的最大平滑连通仿射子群。则 \(u\) 经由 \(G\) 的每个极大环面 \(T\) 分解。特别地，\(u\) 经由上面定义的子函子 \(\mathcal Z\) 分解。}
\SGARefLink{sga3:ix:theorem:5:1:bis}{IX~5.1 bis} 的一个简单变体（用包含代替相等）将断言归约到 \(S\) 是代数闭域的谱的情形：先归约到仿射 Noether \(S\)，再到 Artin \(S\)，并使用 \SGARefLink{sga3:ix:theorem:3:6}{IX~3.6}。由于 \(T\) 包含于 \(G\) 的最大平滑连通仿射子群 \(V\)，可归约到 \(G=V\)，从而归约到 \SGARefLink{sga3:xii:4.5}{4.5} 证明的仿射情形。
\paragraph{\textbf{命题 8.8。}}
\SGARefTarget{sga3:xii:8.8}
\emph{设 \(G\) 是平滑有限呈示且纤维连通的 \(S\)-群预概形。}
\begin{enumerate}
\item[\emph{(a)}] \emph{若 \(S=\operatorname{Spec}k\)，则 \(G\) 存在约化中心 \(Z\)。当 \(G\) 是阿贝尔簇按平滑连通仿射代数群 \(V\) 的扩张时，例如当 \(k\) 代数闭时，\(Z\) 也是 \(V\) 的约化中心，且是 \(V\) 的最大中心乘法型子群。}
\item[\emph{(b)}] \emph{设 \(Z\) 是 \(G\) 的乘法型群子预概形。则 \(Z\) 是 \(G\) 的约化中心，当且仅当对每个 \(s\in S\)，\(Z_s\) 是 \(G_s\) 的约化中心。此时，它是 \(G\) 中满足对每个 \(s\in S\)，\(K_s\) 包含于 \(G_s\) 的约化中心的最大乘法型子群 \(K\)。更一般地，设 \(u:H\to G\) 是同态，其中 \(H\) 是 \(S\) 上有限型乘法型群，并设对 \(S\) 上的每个代数闭域 \(k\)，\(u_k\) 经由 \(G_k\) 的最大平滑连通仿射子群分解。则 \(u\) 经由 \(Z\) 分解，特别地，它是中心的。}
\item[\emph{(c)}] \emph{若 \(G\) 对 fpqc 拓扑局部地存在极大环面，则 \(G\) 存在约化中心。}
\item[\emph{(d)}] \emph{若 \(T\) 是 \(G\) 的极大环面，则}
\[
  T\cap\operatorname{Centr}(G)
  =\ker\bigl(T\longrightarrow\operatorname{GL}(\mathfrak g)\bigr),
\]
\emph{且此群是 \(G\) 的约化中心。}
\item[\emph{(e)}] \emph{设 \(Z\) 是 \(G\) 的约化中心，并设 \(G'=G/Z\) 可表，例如 \(S\) 是 Artin 的时如此。则单位子群是 \(G'\) 的约化中心，且}
\[
  T'\longmapsto T=u^{-1}(T')
\]
\emph{是从 \(G'\) 的极大环面集合到 \(G\) 的极大环面集合的双射。}
\end{enumerate}
\begin{proof}
\emph{\SGAInlineRefTarget{sga3:xii:8.8:item:a:02}(a)} 令 \(S=\operatorname{Spec}k\)。为证明存在性，可取 \(k\) 的代数闭包，从而设 \(G\) 是阿贝尔簇按平滑连通仿射代数群 \(V\) 的扩张。对 \(k\) 上的每个 \(S'\)，\(G_{S'}\) 的极大环面就是 \(V_{S'}\) 的极大环面：这来自 \SGARefLink{sga3:ix:corollary:5:2}{IX~5.2}，以及代数闭域上从环面到阿贝尔簇的同态平凡这一事实。因此，用 \(G\) 定义的函子 \(\mathcal Z\) 与用 \(V\) 定义的函子相同，从而归约到 \(G\) 仿射的情形。
函子 \(\mathcal T\) 可表，且必在 \(k\) 上“本质自由”（\SGARefLink{sga3:viii:definition:6:1}{VIII~6.1}）；故 \(\mathcal T_G\) 在 \(G\) 上本质自由。由于万有环面是 \(\mathcal T_G=G_{\mathcal T}\) 的闭子概形（例如由 \SGARefLink{sga3:viii:corollary:5:7}{VIII~5.7}），\SGARefLink{sga3:viii:theorem:6:4}{VIII~6.4} 表明
\[
  \mathcal Z=
  \prod_{\mathcal T_G/G}
  \widetilde{\mathcal T}/\mathcal T_G
\]
可由 \(G\) 的闭子概形表示。因此它是群子概形。它为乘法型：取代数闭包后，\(G\) 存在极大环面 \(T\)，且由定义 \(Z\subset T\)；乘法型群的子群仍为乘法型（IX~8）。因此 \(Z\) 是约化中心。\SGARefLink{sga3:xii:8.7}{引理 8.7} 表明，它是 \(V\) 的最大中心乘法型子群。
\emph{\SGAInlineRefTarget{sga3:xii:8.8:item:b:02}(b)} 只须证明充分性。设 \(Z\) 为乘法型，且对每个 \(s\in S\)，\(Z_s\) 是 \(G_s\) 的约化中心。\SGARefLink{sga3:xii:8.7}{引理 8.7} 首先表明，\(Z\) 包含于 \(G\) 的每个极大环面，且这在每次基变换后仍成立。只须再说明，若 \(g\) 是 \(G\) 在 \(S\) 上的截面，其到 \(S\) 上的每个 \(S'\) 的拉回都是 \(G_{S'}\) 的每个极大环面的截面，则 \(g\) 是 \(Z\) 的截面。
照例，由于 fpqc 层
\[
  \mathcal Z_0=
  \prod_{\mathcal T_G/G}
  \widetilde{\mathcal T}/\mathcal T_G
\]
与环的滤过直极限交换，可依次归约到 \(S\) 仿射、Noether、Artin 的情形。由 \SGARefLink{sga3:xii:7.1:item:d}{7.1(d)}，\(\mathcal T\) 此时可由 \(S\) 上的平滑预概形表示。仿照\textup{\SGARefLink{sga3:xii:8.8:item:a:02}{(a)}} 论证，\(\mathcal Z_0\) 可由 \(G\) 的闭子概形 \(Z_0\) 表示。已知 \(Z\subset Z_0\)；由假设，在剩余域 \(k\) 上有 \(Z_k=(Z_0)_k\)，且 \(Z\) 在 \(S\) 上平坦。因此 \(Z=Z_0\)，证明了 \(Z\) 是约化中心。\textup{\SGARefLink{sga3:xii:8.8:item:b:02}{(b)}} 的其余断言来自 \SGARefLink{sga3:xii:8.7}{8.7}。
\emph{\SGAInlineRefTarget{sga3:xii:8.8:item:c:02}(c)} 立即来自\textup{\SGARefLink{sga3:xii:8.8:item:d:02}{(d)}}。
\emph{\SGAInlineRefTarget{sga3:xii:8.8:item:d:02}(d)} 这里 \(\mathfrak g\) 表示 \(G\) 的 Lie 代数，而 \(T\to\operatorname{GL}(\mathfrak g)\) 由伴随表示诱导。显然，
\[
  T\cap\operatorname{Centr}(G)
  \subset
  \ker\bigl(T\longrightarrow\operatorname{GL}(\mathfrak g)\bigr).
\]
反向包含的证明与 \SGARefLink{sga3:xii:4.7:item:d}{4.7(d)} 相同。置
\[
  Z=\ker\bigl(T\longrightarrow\operatorname{GL}(\mathfrak g)\bigr).
\]
则 \(Z\) 为乘法型（\SGARefLink{sga3:ix:theorem:6:8}{IX~6.8}）。由\textup{\SGARefLink{sga3:xii:8.8:item:b:02}{(b)}}，可在域上检验它是约化中心。令 \(Z_0\) 为\textup{\SGARefLink{sga3:xii:8.8:item:a:02}{(a)}} 所给的约化中心。显然
\[
  Z_0\subset T\cap\operatorname{Centr}(G)=Z.
\]
反之，\(Z\) 是包含于平滑连通仿射群 \(T\) 的中心乘法型子群，所以\textup{\SGARefLink{sga3:xii:8.8:item:a:02}{(a)}} 给出 \(Z\subset Z_0\)。因此 \(Z=Z_0\)。
\emph{\SGAInlineRefTarget{sga3:xii:8.8:item:e:02}(e)} 在 \SGARefLink{sga3:xii:7.1:item:f}{7.1(f)} 的假设下，置 \(Z=\ker u\)，并设对 \(S\) 上的每个代数闭域 \(k\)，\(Z_k\) 包含于 \(G_k\) 的一个极大环面。容易验证
\[
  T'\longmapsto u^{-1}(T')
\]
是从 \(G'\) 的极大环面到 \(G\) 的极大环面的双射。将此观察用于\textup{\SGARefLink{sga3:xii:8.8:item:e:02}{(e)}} 的情形，便证明了断言。
\end{proof}
\paragraph{\textbf{注 8.9。}}
\SGARefTarget{sga3:xii:8.9}
\begin{enumerate}
\item[\emph{(a)}] \SGARefLink{sga3:xii:8.8}{8.8} 的证明不依赖第 4 节的结果；特别地，\SGARefLink{sga3:xii:8.8:item:a:02}{8.8(a)} 的证明不使用证明稍显繁琐的 \SGARefLink{sga3:xii:4.4}{4.4}。
\item[\emph{(b)}] 无论可表与否，\SGARefLink{sga3:xii:8.6}{8.6} 所考虑的子函子 \(\mathcal Z\) 总是中心的。\ 因此，在所有情形中都将它称为 \(G\) 的约化中心，或许是自然的。
\item[\emph{(c)}] 同样可推广 \SGARefLink{sga3:xii:4.9}{4.9}。若 \(G\) 是代数闭域上的平滑连通代数群，则 \(G\) 是阿贝尔簇按平滑连通幺幂代数群的扩张——由 \SGARefLink{sga3:xii:8.3}{8.3}，这等价于 \(G\) 的约化秩为零——当且仅当 \(G\) 的约化中心为单位群，且 \(G\) 的 Lie 代数幂零。
\end{enumerate}
\section{补充：群概形的作用与不动点}
\SGARefTarget{sga3:xii:section:9}
本节于 2008 年 1 月增补，旨在研究群概形作用下的不动点。
\subsection{不动点函子的可表性}
\SGARefTarget{sga3:xii:9.1}
设 \(S\) 是概形，\(G\) 是 \(S\)-群概形，并设 \(G\) 作用于 \(S\)-概形 \(X\)。定义 \(X\) 在 \(G\) 下的不动点子函子 \(X^G\)：对每个 \(S\)-概形 \(T\)，要求 \(X^G(T)\) 为 \(X(T)\) 中在 \(G\) 下不动的点组成的子集（VIII~6\SGARefLink{sga3:xii:8.8:item:e:02}{(e)}）。
回顾 Raynaud--Gruson，\S3.3 及 Raynaud 所引入的 \(S\)-纯概形的概念。设 \(G\) 是平坦有限呈示 \(S\)-群概形。由于 \(G\) 的纤维没有嵌入分支，纯性具有如下形式：若对每个闭点为 \(t\) 的严格 Hensel 局部 \(S\)-概形 \(T\)，\(G\times_S T\) 的某个纤维的每个一般点 \(\widetilde x\) 都特殊化为 \(G_t\) 的一点，或等价地，\(\widetilde x\) 在 \(G\times_S T\) 中的闭包与 \(G_t\) 相交，则称 \(G\) 是 \emph{\(S\)-纯的}。
\paragraph{\textbf{例。}}
\begin{enumerate}
\item[\textup{(1)}] 若 \(G\) 在 \(S\) 上拟有限且分离，则 \(G\) 是 \(S\)-纯的，当且仅当它在 \(S\) 上有限。
\item[\textup{(2)}] 若 \(S=\operatorname{Spec}R\) 且 \(G=\operatorname{Spec}A\)，则 \(G\) 是 \(S\)-纯的，当且仅当 \(A\) 是投射 \(R\)-模（Raynaud--Gruson，3.3.5）。特别地，可对角化的 \(G\) 是 \(S\)-纯的。
\item[\textup{(3)}] 每个乘法型群概形都是 \(S\)-纯的，因为纯性对 \(S\) 上的平展拓扑是局部的。
\item[\textup{(4)}] 若 \(G\) 在 \(S\) 上固有，或其纤维不可约，或 \(S\) 是半局部 Artin 的，则概形 \(G\) 是 \(S\)-纯的。
\item[\textup{(5)}] 特别地，每个约化 \(S\)-群概形都是 \(S\)-纯的。
\end{enumerate}
\paragraph{\textbf{命题 9.2。}}
\SGARefTarget{sga3:xii:9.2}
\emph{设 \(G\) 是在 \(S\) 上有限呈示、平坦且 \(S\)-纯的 \(S\)-群概形，作用于有限呈示 \(S\)-概形 \(X\)。则不动点函子 \(X^G\) 可由 \(X\) 的一个在 \(S\) 上有限呈示的闭子概形表示。}
事实上，设 \(a\in X(S)\)，\(H\) 是 \(G\) 中固定 \(a\) 的群子概形。点 \(a\) 在 \(G\) 下不动，当且仅当 \(H=G\)。\(H\) 与 \(G\) 的相合函子——它在 \(S\)-概形 \(T\) 处的值，当
\[
  H\times_S T\longrightarrow G\times_S T
\]
为双射时是独点集，否则为空集——可由 \(S\) 的某个闭子概形表示，且该闭子概形由有限型理想层定义（Raynaud--Gruson，4.1.1）。取 \(S=X\)，且取 \(a\) 为 \(X\) 的万有点，即
\[
  X(X)=\operatorname{Hom}_S(X,X).
\]
的恒等元素。
\setcounter{subsection}{2}
\subsection{无穷小障碍}
\SGARefTarget{sga3:xii:9.3}
设 \(G\) 是平坦有限呈示 \(S\)-群概形，在平滑 \(S\)-预概形 \(X\) 上左作用。我们研究 \(X^G\) 的形式平滑性。
设 \(S_0\) 是 \(S\) 的闭子概形，由满足 \(\mathcal I^2=0\) 的拟凝聚理想层 \(\mathcal I\) 定义，记
\[
  j:S_0\longrightarrow S,\qquad
  G_0=G\times_S S_0,\qquad
  X_0=X\times_S S_0.
\]
设 \(\epsilon_0\in X^G(S_0)\) 是可提升到 \(X(S)\) 的点。我们研究将 \(\epsilon_0\) 提升到 \(X^G(S)\) 的障碍。
由于 \(X\) 平滑，\(\epsilon_0\) 在 \(X(S)\) 中的提升所成集合，是阿贝尔群
\[
  \operatorname{Hom}_{\mathcal O_{S_0}}
  \bigl(\epsilon_0^*\Omega^1_{X_0/S_0},\mathcal I\bigr)
\]
下的平凡主齐性空间（\SGARefLink{sga3:iii:proposition:0-2}{III~0.2}、\SGARefLink{sga3:iii:remark:0-3}{0.3}）。置
\[
  L_0=
  \operatorname{Hom}_{\mathcal O_{S_0}}
  \bigl(\epsilon_0^*\Omega^1_{X_0/S_0},\mathcal I\bigr).
\]
这是 \(\mathcal O_{S_0}\)-模。由于 \(\epsilon_0\) 在 \(G\) 下不动，它自然是 \(G_0\)-\(\mathcal O_{S_0}\)-模。令
\[
  \rho_0:G_0\longrightarrow\operatorname{Aut}_{S_0}(L_0),
  \qquad
  \rho:G\longrightarrow\operatorname{Aut}_S(j_*L_0)
\]
为相应表示。
\paragraph{\textbf{引理 9.4。}}
\SGARefTarget{sga3:xii:9.4}
\emph{存在一个典范定义的类}
\[
  c(\epsilon_0)\in H^1(G,j_*L_0)
  \simeq H^1(G_0,L_0)
\]
\emph{使得 \(\epsilon_0\) 可提升到 \(X^G(S)\)，当且仅当 \(c(\epsilon_0)=0\)。}
所列伴随对应就是 \SGARefLink{sga3:iii:lemma:1-1-2}{III~1.1.2} 中的对应。我们在 \(S\) 的小 fppf 景上工作。对每个平坦有限呈示 \(S\)-概形 \(T\)，记
\[
  G_T=G\times_S T,\qquad X_T=X\times_S T.
\]
考虑此景上如下定义的层 \(\mathcal A\)
\[
  \mathcal A(T)=
  \{\,\epsilon_0\times_S T\text{ 在 }X(T)\text{ 中的提升}\,\}.
\]
由于 \(L_0\) 的构造与平坦基变换交换，\(\mathcal A(T)\) 是
\[
  H^0(T,j_*L_0).
\]
下的平凡主齐性空间。此外，由于 \(\epsilon_0\) 不动，每个 \(g\in G(T)\) 通过仿射自同构作用于 \(\mathcal A(T)\)，且与 \(G\) 在 \(j_*L_0\) 上的作用相容：
\[
  g(a_T+v)=g(a_T)+\rho(g)(v),
  \qquad v\in H^0(T,j_*L_0).
\]
由于 \(G\) 平坦且有限型，可取 \(T=G\)，且取 \(g\) 为万有点 \(\operatorname{id}_G\)。这样得到 fppf 层 \(G\) 在 \(\mathcal A\) 上的作用。
设 \(a\) 为 \(\epsilon_0\) 在 \(X(S)\) 中的一个提升，用 \(g^\sharp\) 表示 \(G\) 的万有点。定义
\[
  v^\sharp\in H^0(G,j_*L_0)
\]
其定义关系为
\[
  \rho(g^\sharp)(v^\sharp)=g^\sharp\cdot a_G-a_G.
\]
对每个 \(S\)-预概形 \(Y\)，置
\[
  z(Y):G(Y)\longrightarrow H^0(Y,j_*L_0),
  \qquad g\longmapsto v^\sharp(g).
\]
这定义了一个 \(1\)-余循环 \(z\in Z^1(G,j_*L_0)\)（第 I 篇，\S5）。其类 \(c(\epsilon_0)\) 不依赖于 \(a\) 的选择。特别地，若 \(\epsilon_0\) 可提升到 \(X^G(S)\)，则 \(c(\epsilon_0)=0\)。
反之，若 \(c(\epsilon_0)=0\)，存在截面
\[
  w\in H^0(S,j_*L_0)
\]
使对每个 \(S\)-预概形 \(Y\) 及每个 \(g\in G(Y)\)，
\[
  z(Y)(g)=g\cdot w_Y-w_Y.
\]
将其用于 \(Y=G\)、\(g=g^\sharp\)，可得 \(a-w\in X^G(S)\)。因此，\(\epsilon_0\) 可提升到 \(X^G(S)\)，当且仅当 \(c(\epsilon_0)=0\)。
\paragraph{\textbf{注 9.5。}}
\SGARefTarget{sga3:xii:9.5}
当 \(S\) 仿射，且 \(G\) 在 \(S\) 上仿射、平坦、有限型时，此障碍可细化为类
\[
  \widetilde c(\epsilon_0)\in H^1_{G_0}(X_0,L_0),
\]
其中右边为 Wewers 定义的 \(G_0\)-等变上同调群（附录 C）。此时不必假设 \(\epsilon_0\) 可提升到 \(X(S)\)。
\setcounter{subsection}{5}
\subsection{不动点的平滑性}
\SGARefTarget{sga3:xii:9.6}
\paragraph{\textbf{定理 9.7。}}
\SGARefTarget{sga3:xii:9.7}
\emph{设 \(G\) 是 Noether 基 \(S\) 上的平坦有限型 \(S\)-群概形，作用于平滑 \(S\)-概形 \(X\)。设 \(G\) 是 \(S\)-纯的，且对 \(S\) 的每个剩余域代数闭的几何点 \(s\)，}
\[
  H^1\bigl(G_s,(\Omega^1_{X_s})^*\bigr)=0.
\]
\emph{则不动点函子 \(X^G\) 可由 \(X\) 的一个在 \(S\) 上平滑的闭子概形表示。}
由 9.2，\(X^G\) 可由 \(X\) 的一个在 \(S\) 上有限呈示的闭子概形表示。可设 \(S=\operatorname{Spec}A\)，其中 \(A\) 局部。由 SGA~1，III~3.1(iii bis) 的平滑性判据，只须对
\[
  S'=\operatorname{Spec}A'
\]
及其闭子概形
\[
  S''=\operatorname{Spec}(A'/I'),
\]
验证形式平滑性，其中 \(A'\) 是 Artin 局部环，\(I'\) 是平方为零的理想。通过逐层归约，可设 \(I'\) 被 \(A'\) 的极大理想零化。若 \(k\) 是 \(A'\) 的剩余域，则 \(I'\) 因而是 \(k\)-向量空间。
设 \(\epsilon''\in X^G(S'')\) 是要提升到 \(X^G(S')\) 的点，置
\[
\begin{gathered}
  X'=X\times_S S',\qquad G'=G\times_S S',\\
  X''=X\times_S S'',\qquad G''=G\times_S S''.
\end{gathered}
\]
由于 \(S'\) 仿射，且 \(X'\) 在 \(S'\) 上平滑，\(\epsilon''\) 可提升到 \(X(S')\)。由\SGARefLink{sga3:xii:9.4}{引理 9.4}，将 \(\epsilon''\) 提升到 \(X^G(S')\) 的障碍是类
\[
  c(\epsilon'')\in H^1(G'',L''),
\]
其中
\[
  L''=
  \operatorname{Hom}_{\mathcal O_{S''}}
  \bigl((\epsilon'')^*\Omega^1_{X''/S''},I'\bigr).
\]
由于 \(L''\) 是 \(k\)-向量空间，有典范同构
\[
  H^1(G'',L'')
  \simeq
  H^1(G''\times_{A'/I'}k,L'').
\]
设 \(\bar k\) 是 \(k\) 的代数闭包。在仿射情形由 \SGARefLink{sga3:ix:theorem:3:1}{IX~3.1}，在一般情形由\SGARefLink{sga3:xii:9.11}{引理 9.11}，有同构
\[
  H^1(G''\times_{A'/I'}k,L'')\otimes_k\bar k
  \xrightarrow{\ \sim\ }
  H^1(G''\times_{A'/I'}\bar k,L''\otimes_k\bar k).
\]
表示 \(L''\otimes_k\bar k\) 是
\[
  (\Omega^1_{X_{\bar k}})^*.
\]
的若干副本的直和。因此右边的群由假设为零。可知 \(H^1(G'',L'')=0\)，故障碍消失，\(X^G\) 形式平滑。
\paragraph{\textbf{推论 9.8。}}
\SGARefTarget{sga3:xii:9.8}
\emph{设 \(G\) 是平坦有限呈示 \(S\)-群概形，作用于平滑有限呈示 \(S\)-概形 \(X\)。假设 \(G\) 有一个合成列，其因子属于以下类型之一：}
\begin{enumerate}
\item[\emph{(1)}] \emph{阿贝尔 \(S\)-概形，即纤维为阿贝尔簇的平滑 \(S\)-群概形；}
\item[\emph{(2)}] \emph{乘法型 \(S\)-群概形；}
\item[\emph{(3)}] \emph{次数在 \(S\) 上可逆的有限平展 \(S\)-群概形；}
\item[\emph{(4)}] \emph{当 \(S\) 是 \(\mathbb Q\)-概形时，约化 \(S\)-群概形。}
\end{enumerate}
\emph{则不动点函子 \(X^G\) 可由 \(X\) 的一个在 \(S\) 上有限呈示且平滑的闭子概形表示。}
事实上，设
\[
  1\longrightarrow G'\longrightarrow G\longrightarrow G''
  \longrightarrow 1
\]
是满足假设的 \(S\)-群正合列。群 \(G''\) 作用于 \(X^{G'}\)，而 \(X^G\) 是 \(X^{G'}\) 在 \(G''\) 下的不动点函子。因此归约到群属于上述四种类型之一的情形。通过取极限，可设 \(S\) 是 Noether 的。对几何纤维 \(G_s\) 及 \(G_s\) 的有限维线性表示 \(V\)，验证\SGARefLink{sga3:xii:9.7}{定理 9.7} 的上同调判据。
对阿贝尔概形，从 \(G_s^n\) 到 \(V\) 的每个态射都为常值；因此
\[
  H^i(G_s,V)=0\qquad(i>0).
\]
对乘法型群概形，\(G_s\) 可对角化，\SGARefLink{sga3:i:theorem:5-3-3}{I~5.3.3} 再次给出
\[
  H^i(G_s,V)=0\qquad(i>0).
\]
在剩下两种情形中，断言来自 \(G_s\) 的线性表示的半单性。若 \(G_s\) 有限平展，且次数在 \(\kappa(s)\) 中可逆，则它是阶可逆的常值有限群，且 \(G_s\) 的每个线性表示都半单。在约化情形，\(\kappa(s)\) 是特征零代数闭域，\(G_s\) 的每个线性表示都半单（Tauvel--Yu，定理 27.3.3）。
\paragraph{\textbf{注 9.9。}}
\SGARefTarget{sga3:xii:9.9}
这特别适用于 \(G/S\) 在平滑 \(S\)-群概形 \(H\) 上的作用。若 \(G\) 为乘法型，通过内自同构作用于 \(H\)，其中 \(H\) 在 \(S\) 上平滑且仿射，就重新得到 \SGARefLink{sga3:xi:corollary:5:3}{XI~5.3}：\(G\) 在 \(H\) 中的中心化子平滑。
\paragraph{\textbf{注 9.10。}}
\SGARefTarget{sga3:xii:9.10}
当 \(S=\operatorname{Spec}k\)，且 \(G\) 是代数闭域 \(k\) 上的线性约化代数群时，此结果由 Fogarty（定理 5.4）和 Iversen（命题 1.3）独立得到。
\SGARefLink{sga3:xii:9.7}{定理 9.7} 的证明使用了以下引理，对仿射群概形而言它是熟知的（\SGARefLink{sga3:ix:theorem:3:1}{IX~3.1}）。
\paragraph{\textbf{引理 9.11。}}
\SGARefTarget{sga3:xii:9.11}
\emph{设 \(G\) 是仿射概形 \(S=\operatorname{Spec}A\) 上的群概形，\(f:S'=\operatorname{Spec}A'\to S\) 平坦。置 \(G'=G\times_S S'\)。则对每个拟凝聚 \(G\)-\(\mathcal O_S\)-模 \(M\)，有典范同构}
\[
  H^i(G,M)\otimes_A A'
  \xrightarrow{\ \sim\ }
  H^i(G',M\otimes_A A')
  \qquad(i\geq 0).
\]
对 \(n\geq1\)，\(n\) 次 Hochschild 余链由其在
\[
  (\operatorname{id},\ldots,\operatorname{id})
  \in G(G)\times\cdots\times G(G).
\]
处的值确定。因此有 \(A\)-模同构
\[
  C^n(G,M)
  \simeq
  \Gamma\bigl(G^n,M\otimes_A\mathcal O_{G^n}\bigr).
\]
群 \(H^n(G,M)\) 是复形
\[
  M\longrightarrow
  \Gamma(G,M\otimes_A\mathcal O_G)\longrightarrow
  \Gamma(G^2,M\otimes_A\mathcal O_{G^2})\longrightarrow
  \Gamma(G^3,M\otimes_A\mathcal O_{G^3})\longrightarrow\cdots.
\]
的第 \(n\) 个上同调群。由于 \(A'\) 在 \(A\) 上平坦，EGA IV\(_1\)，1.7.21 给出
\[
\begin{split}
  \Gamma\bigl(G^n,M\otimes_A\mathcal O_{G^n}\bigr)\otimes_A A'
  \xrightarrow{\ \sim\ }{}&
  \Gamma\bigl((G')^n,\\
  &(M\otimes_A A')\otimes_{A'}
  \mathcal O_{(G')^n}\bigr).
\end{split}
\]
取上同调即证断言。
\begin{thebibliography}{Ray72}
\SGARefTarget{sga3:xii:bibliography}
\addcontentsline{toc}{section}{参考文献}
\item[]\emph{编者注：以下是本篇引用的补充参考文献。}
\bibitem[Fo73]{sga3-xii-Fo73}
J.\ Fogarty，
\emph{不动点概形〔Fixed point schemes〕}，
Amer.\ J.\ Math.\ \textbf{95} (1973)，35--51。
\bibitem[RG71]{sga3-xii-RG71}
M.\ Raynaud 和 L.\ Gruson，
\emph{平坦性与投射性判据〔Critères de platitude et de projectivité〕}，
Invent.\ Math.\ \textbf{13} (1971)，1--89。
\bibitem[Iv72]{sga3-xii-Iv72}
B.\ Iversen，
\emph{环面作用于代数簇的不动点公式〔A fixed point formula for action of tori on algebraic varieties〕}，
Invent.\ Math.\ \textbf{16} (1972)，229--236。
\bibitem[Ray72]{sga3-xii-Ray72}
M.\ Raynaud，
\emph{代数几何中的平坦模〔Flat modules in algebraic geometry〕}，
Compositio Math.\ \textbf{24} (1972)，11--31。
\bibitem[TY06]{sga3-xii-TY06}
P.\ Tauvel 和 R.\ W.\ T.\ Yu，
\emph{Lie 代数与代数群〔Lie algebras and algebraic groups〕}，
Springer-Verlag，2006。
\bibitem[We05]{sga3-xii-We05}
S.\ Wewers，
\emph{带群概形作用的曲线的形式形变〔Formal deformation of curves with group scheme action〕}，
Ann.\ Inst.\ Fourier (Grenoble) \textbf{55} (2005)，1105--1165。
\end{thebibliography}
